\begin{theorem}[Two-gather round refinement]
  \label{thm:gbca-round-refines-specification}
  \lean{PLTS.ABA.GBCA.count_aboveThreshold_of_subMap}\lean{PLTS.ABA.GBCA.boundOfCore_of_aboveThreshold}\lean{PLTS.ABA.GBCA.count_boundOfCore_belowThreshold}\lean{PLTS.ABA.GBCA.ByAFW.Invariant}\lean{PLTS.ABA.GBCA.ByAFW.ExclusionWitness}\lean{PLTS.ABA.GBCA.ByAFW.SpecificationRelation}\lean{PLTS.ABA.GBCA.ByAFW.refinesSpecification}\leanok
  \uses{def:aba-round-over-gather-specifications, def:aba-gbca-spec}
  The two-gather round is forward-simulated by the GBCA specification instance (Transition System~2): \leandecl{PLTS.ABA.GBCA.ByAFW.refinesSpecification}.
  The exclusion set and the grade guard are witnessed on the two cores, each written once, --- a bit is excluded when the first core carries it on fewer than $|S| - f$ entries --- and excludes fire on demand under the returns that need them.
  The value of a return of grade 1 or 2 is the round's bound bit, so the specification's guard on its announcement and its guard pair on the value are the same exclusion, and the round's announcement is matched by the specification's.
  The candidate's witness is recorded at the first gather's return and consumed at the second gather's call, and the grade's at the second gather's return and consumed at the graded return, the program's variables carrying each across its two events.
\end{theorem}

\begin{proof}\leanok
  \uses{def:aba-round-over-gather-specifications, def:aba-gbca-spec, def:forwardsim}
  See the Lean proof of \leandecl{PLTS.ABA.GBCA.ByAFW.refinesSpecification}.
\end{proof}
